Cubic completion and aperiodic transport specify two senses
of preservation articulated in the development. The former preserves
unity under projection of normalized measures and distinguishes the
initial cube, the punctured completion of \(999\) elements, and restoration
of the apex. The latter preserves the trajectory record during
phase return. Monodromy and the mathematical model of an aperiodic
time crystal express this second mechanism through maps and
observables, not merely through a helical image. The tritic
exponential, for its part, unifies the elliptic, parabolic, and
hyperbolic realizations without suppressing their respective quadratic relations.

The reduced Planck constant, \(\hbar\), occupies an explicit position
in this architecture. The order-\(16\) determinantal quotient acts
in the norm determining the decade \(-34\); the character \(H_5\) and its
regional continuation determine the two action sections whose
quotient enters the electronic coordinate. Comparing the
preceding section with \(h/(2\pi)\) shows a relative difference of
order \(10^{-15}\), and the returned section preserves the compensation
used by the electron. The action scale thus links the determinantal norm
and regional character to the electronic quotient: the compensation belongs
to the exact construction, whereas the proximity of the preceding section
to \(h/(2\pi)\) is a numerical comparison, not an exact equality.
